\subsection{Uniform Flexible-Start Study}
\label{subsec:uniform_results}
The new comparison uses the common policy in \Cref{subsec:uniform_flexible}.
Both encodings use CPLEX 12.6.0.0, one thread, presolve on, and a 180-second limit.
They share the hour engine; the comparison measures maintenance encoding, not independent physical models.
All loaded incumbents, including time-limited ones, are checked by chronological replay.
An additional audit reconstructs cost, hours, timing, routes, actual start-day counts and bay occupancy
directly from the saved decisions and input JSON, without the model's state values.
\begin{table}[t]\centering\scriptsize
\caption{Uniform flexible-start encodings. Wall time is measured around the solver call.
Status is the interface termination code; a time-limited incumbent is not an optimum.}
\label{tab:uniform_flexible}\resizebox{\linewidth}{!}{%
\begin{tabular}{rrl lrrrrl}\toprule
$|P|$ & $H$ & Encoding & Status & Objective & Wall (s) & Vars & Rows & Replay \\\midrule
5 & 30 & Day & optimal & 37583 & 0.84 & 1548 & 5283 & PASS \\
5 & 30 & Event & optimal & 37583 & 0.96 & 1733 & 5468 & PASS \\
5 & 48 & Day & optimal & 40707 & 0.73 & 1560 & 5360 & PASS \\
5 & 48 & Event & optimal & 40707 & 0.70 & 1760 & 5560 & PASS \\
10 & 30 & Day & optimal & 74150 & 51.38 & 6188 & 21220 & PASS \\
10 & 30 & Event & optimal & 74150 & 27.17 & 6918 & 21950 & PASS \\
10 & 49 & Day & optimal & 81391 & 3.15 & 6463 & 21959 & PASS \\
10 & 49 & Event & optimal & 81391 & 3.00 & 7263 & 22759 & PASS \\
15 & 30 & Day & optimal & 110749 & 157.11 & 13900 & 49043 & PASS \\
15 & 30 & Event & maxIterations & 110749 & 182.39 & 15535 & 50678 & PASS \\
15 & 49 & Day & maxIterations & 122169 & 182.57 & 14607 & 50211 & PASS \\
15 & 49 & Event & maxIterations & 122169 & 182.41 & 16407 & 52011 & PASS \\
\bottomrule\end{tabular}}\end{table}
Of 12 runs, 9 terminate optimally and 12 return loaded incumbents; all 12 pass both replay checks.
Unavailable solver bounds and node counts are left blank in raw data, not inferred from incumbent values.
The timings are descriptive workstation measurements, not isolated repeated-runtime trials or evidence of general speedup.
\begin{table}[t]\centering\scriptsize
\caption{Repaired day-boundary ablation, including the first-day opening row. This is not the arbitrary flexible-start model.}
\label{tab:repaired_prefix}\resizebox{\linewidth}{!}{%
\begin{tabular}{rr rrr r ll}\toprule
$|P|$ & $H$ & State LP & Interval LP & Prefix LP & Common MIP & Interval vars/rows & Prefix vars/rows \\\midrule
10 & 30 & 72050.00 & 67850.11 & 67850.11 & 74150 & 3120/12395 & 3440/12705 \\
10 & 49 & 79057.67 & 74391.13 & 74391.13 & 81391 & 3900/20798 & 4410/21298 \\
5 & 30 & 36516.33 & 34383.06 & 34383.06 & 37583 & 1020/4498 & 1180/4653 \\
5 & 48 & 39557.00 & 37207.06 & 37207.06 & 40707 & 1335/8625 & 1585/8870 \\
\bottomrule\end{tabular}}\end{table}
All 12 repaired LP and all 12 integer runs terminate optimally. Saved integer schedules pass the day-boundary hour replay.
The interval/prefix LP equality follows from telescoping with identical bounds and opening domains; equality to the stronger reset-state LP is not claimed.
The raw snapshots are \path{repaired_hour_prefix_lp_20261021.csv} and \path{repaired_hour_prefix_milp_20261021.csv}.
\begin{figure}[t]\centering
\includegraphics[width=0.86\linewidth]{figures/uniform_variables_h30_20261021.png}
\caption{Uniform flexible model size at observed $H=30$, including the P=15 models.}
\end{figure}
\begin{figure}[t]\centering
\includegraphics[width=0.86\linewidth]{figures/uniform_wall_h30_20261021.png}
\caption{Uniform flexible solve time at $H=30$ on common optimal cells only (P=5 and P=10).}
\end{figure}
\begin{table}[t]\centering\scriptsize
\caption{Fixed 22:00 restriction on the four P=5/10 cases, with the same 180-second limit.}
\label{tab:uniform_night}\begin{tabular}{rrllrrl}\toprule
$|P|$ & $H$ & Encoding & Status & Objective & Wall (s) & Replay \\\midrule
5 & 30 & Day & optimal & 37583 & 0.78 & PASS \\
5 & 30 & Event & optimal & 37583 & 0.78 & PASS \\
5 & 48 & Day & optimal & 40707 & 0.69 & PASS \\
5 & 48 & Event & optimal & 40707 & 0.69 & PASS \\
10 & 30 & Day & optimal & 74150 & 12.76 & PASS \\
10 & 30 & Event & optimal & 74150 & 12.75 & PASS \\
10 & 49 & Day & optimal & 81391 & 2.15 & PASS \\
10 & 49 & Event & optimal & 81391 & 2.18 & PASS \\
\bottomrule\end{tabular}\end{table}
The nightly model is a restricted flexible model, not Khaled's daily-quota/calendar-spacing model.
\begin{table}[t]\centering\scriptsize
\caption{Uniform flexible three-seed subset: wall time as mean $\pm$ sample standard deviation, optimal runs only.}
\label{tab:uniform_seeds}\begin{tabular}{rrllr}\toprule
$|P|$ & $H$ & Encoding & Optimal / runs & Wall (s) \\\midrule
5 & 30 & Day & 3/3 & $0.874\pm0.027$ \\
5 & 30 & Event & 3/3 & $0.892\pm0.084$ \\
5 & 48 & Day & 3/3 & $0.714\pm0.005$ \\
5 & 48 & Event & 3/3 & $0.706\pm0.009$ \\
10 & 30 & Day & 3/3 & $72.943\pm52.382$ \\
10 & 30 & Event & 3/3 & $22.516\pm5.142$ \\
10 & 49 & Day & 3/3 & $2.705\pm0.072$ \\
10 & 49 & Event & 3/3 & $2.893\pm0.076$ \\
\bottomrule\end{tabular}\end{table}
Three seeds test limited instance variation; seed zero overlaps the six-case study and is not an additional independent sample.
Raw new results use the \path{uniform_flexible_main_20261021.csv},
\path{uniform_night_subset_20261021.csv} and \path{uniform_flexible_multiseed_20261021.csv}
snapshots in \path{results/tables/}, each with saved schedules, source archive and hashed metadata.
The following 2026-10-20 tables are historical snapshots with different policies; their values are not results for these uniform encodings.
